\begin{figure}[htbp!]
	\centering
	\begin{tikzpicture}[
		x=1.5cm, y=1.5cm, font=\scriptsize,
		trave/.style={line width=2.0pt, line cap=round, line join=round, black},
		forza/.style={-{Stealth[length=2mm]}, line width=0.9pt, black},
		carico/.style={-{Stealth[length=1.6mm]}, line width=0.6pt, black},
		bordo/.style={line width=0.5pt, black},
		nodoV/.style={circle, fill=white, draw=black, inner sep=0pt,
			minimum size=5pt, line width=0.8pt},
		nodo/.style={circle, fill=black, inner sep=0pt, minimum size=3pt},
		blob/.style={draw=gray!45, fill=gray!10, line width=0.7pt, rounded corners=10pt},
		]
		\newcommand{\fin}[3]{\pgfmathsetmacro{\bk}{#2+180}%
			\draw[forza] ($(#1)+(\bk:0.46)$)--(#1);
			\node[font=\scriptsize] at ($(#1)+(\bk:0.74)$) {#3};}
		\newcommand{\fout}[3]{%
			\draw[forza] (#1)--($(#1)+(#2:0.46)$);
			\node[font=\scriptsize] at ($(#1)+(#2:0.76)$) {#3};}
		% ===== corpo ADE (alto-sinistra) =====
		\begin{scope}[shift={(0,3.2)}]
			\draw[blob] (-0.92,-0.92) rectangle (1.50,1.92);
			\draw[trave] (0,0)--(0,1)--(0.5,1);
			\foreach \x in {0,0.1,...,0.6}\draw[carico] (\x,0.78)--(\x,1.0);
			\draw[bordo] (0,0.78)--(0.5,0.78);  \node[left] at (-0.04,0.86) {$p$};
			\fin{0,0}{0}{$\tfrac58 p\ell$}   \fin{0,0}{90}{$p\ell$}
			\fin{0.5,1}{180}{$\tfrac58 p\ell$} \fin{0.5,1}{270}{$\tfrac32 p\ell$}
			\node[nodoV] at (0,0){}; \node[nodoV] at (0.5,1){}; \node[nodo] at (0,1){};
			\node at (0.20,0.0){$A$}; \node at (-0.16,1.10){$D$}; \node at (0.68,1.15){$E$};
			%\node[font=\footnotesize] at (0.25,-1.22){corpo $ADE$};
		\end{scope}
		% ===== corpo EFGB (alto-centro) =====
		\begin{scope}[shift={(3.0,3.2)}]
			\draw[blob] (-0.92,-1.00) rectangle (2.15,1.92);
			\draw[trave] (0,1)--(0.5,1)--(0.5,0); \draw[trave] (0.5,1)--(1.5,1);
			\foreach \x in {0,0.1,...,0.5}\draw[carico] (\x,0.78)--(\x,1.0);
			\draw[bordo] (0,0.78)--(0.5,0.78);  \node at (0.25,0.62){$p$};
			\fin{0,1}{0}{$\tfrac58 p\ell$}  \fin{0,1}{90}{$\tfrac32 p\ell$}
			\fin{0.5,0}{180}{$\tfrac58 p\ell$} \fout{0.5,0}{270}{$\tfrac72 p\ell$}
			\fin{1.5,1}{90}{$\tfrac32 p\ell$}
			\node[nodoV] at (0,1){}; \node[nodoV] at (0.5,0){}; \node[nodoV] at (1.5,1){};
			\node[nodo] at (0.5,1){};
			\node at (-0.16,1.14){$E$}; \node at (0.42,1.16){$F$};
			\node at (0.32,0.16){$B$}; \node at (1.64,1.12){$G$};
			%\node[font=\footnotesize] at (0.7,-1.22){corpo $EFGB$};
		\end{scope}
		% ===== corpo CG (alto-destra) =====
		\begin{scope}[shift={(6.5,3.2)}]
			\draw[blob] (-0.85,-0.92) rectangle (0.85,1.92);
			\draw[trave] (0,0)--(0,1);
			\fout{0,0}{270}{$\tfrac12 p\ell$} \fout{0,1}{90}{$\tfrac12 p\ell$}
			\node[nodoV] at (0,0){}; \node[nodoV] at (0,1){};
			\node at (0.20,0.04){$C$}; \node at (0.20,1.0){$G$};
			%\node[font=\footnotesize] at (0,-1.22){corpo $CG$};
		\end{scope}
		% ===== corpo AB (basso-sinistra) =====
		\begin{scope}[shift={(0,1.0)}]
			\draw[blob] (-0.92,-1.00) rectangle (1.94,0.95);
			\draw[trave] (0,0)--(1,0);
			\draw[black,line width=1.0pt,-{Stealth[length=2mm]}]
			(0.5,0)++(215:0.16) arc[start angle=215,end angle=-35,radius=0.16];
			\node at (0.5,0.30){$p\ell^2$};
			\fout{0,0}{180}{$\tfrac58 p\ell$} \fout{0,0}{270}{$p\ell$}
			\fout{1,0}{0}{$\tfrac58 p\ell$}   \fout{1,0}{90}{$p\ell$}
			\node[nodoV] at (0,0){}; \node[nodoV] at (1,0){};
			\node at (-0.02,0.22){$A$}; \node at (1.10,-0.20){$B$};
			%\node[font=\footnotesize] at (0.5,-1.22){corpo $AB$};
		\end{scope}
		% ===== corpo BC (basso-centro) =====
		\begin{scope}[shift={(3.0,1.0)}]
			\draw[blob] (-0.85,-1.00) rectangle (1.80,0.85);
			\draw[trave] (0,0)--(1,0);
			\foreach \x in {0,0.1,...,1.10}\draw[carico] (\x,0.22)--(\x,0.0);
			\draw[bordo] (0,0.22)--(1.0,0.22);  \node[right] at (1.04,0.30){$p$};
			\fin{0,0}{90}{$\tfrac12 p\ell$} \fin{1,0}{90}{$\tfrac12 p\ell$}
			\node[nodoV] at (0,0){}; \node[nodoV] at (1,0){};
			\node at (-0.18,-0.04){$B$}; \node at (1.18,-0.04){$C$};
			%\node[font=\footnotesize] at (0.5,-1.22){corpo $BC$};
		\end{scope}
	\end{tikzpicture}
	\caption{Anelli a tre cerniere: esploso dei cinque corpi
		liberi con i carichi attivi e le forze scambiate alle
		connessioni interne. La spinta~$\tfrac58 p\ell$ si chiude
		nell'anello portato~$ABFD$; la coppia attiva raggiunge
		il suolo attraverso il montante teso~$BF$.}
	\label{fig:anello3c_corpi_liberi}
\end{figure}
